\chapter*{Plan (Brouillon)}

\begin{itemize}
    \item Introduction \textbf{(\textasciitilde 3 p.)}
    \begin{itemize}
        \item "light" > give the entry point on sparse matrices and irregular accesses \textbf{(1-2 p.)}
        \item Mains contributions \textbf{(1 p.)}
        \item introduce the structure of the manuscript
    \end{itemize}
    \item Background \textbf{(\textasciitilde 5 p.)}
    \begin{itemize}
        \item irregular accesses -> def, issues (in memory but not only?), domains/applications (including sparse matrix processing) \textbf{(1 p.)}
        \item Focus on sparse matrices, context: it is widely used, in many domains, sparse, large, structures
        \item Sparse matrices in a kernel: show where the accesses are regular vs irregular (infos on basic sparse formats (other sparse formats in appendix x), tiling, basic kernels, basic algos...) \textbf{(2-3 p.)}
        \item Hardware challenges clearly defined (gather, scatter...) \textbf{(1 p.)}
    \end{itemize}
    \item SoA \textbf{(\textasciitilde 11 p.)}
    \begin{itemize}
        \item Present solutions in SW for CPU/GPU, such as special sparse formats, possibility to play on algo (gust analysis), tiling methods adapted for sparse data (specific sparse format + specific algos) (expliquer clairement que ces solutions sont interressantes et pourquoi on se concentre sur le HW) \textbf{(1-2 p.)}
        \item Present solutions in HW (dedicated accelerators) -> comparison table along common characteristics \textbf{(1-2 p.)}
        \item Presentation of the major accelerators (selection of 3 important: processor assist PHI, standalone GAMMA, multi algo SPADA, ACES) (others in appendix xi) \textbf{(3-4 p.)}
        \item Problem of comparison / comparison of standalone / other analysis (issue of matrix structure) -> Objectives: convince reader that SoA was done in depth, and acknoledge the matrice structure issue \textbf{(1 p.)}
        \item Takeaways -> paths for designing an accelerator (end od SoA paper) \textbf{(1-2 p.)}
        \item We choose a path -> show which one (coarse grain, at the level of the SoA, comparable to SortCache, PHI...)
    \end{itemize}
    \item High level architectural presentation of SpDCache \textbf{(\textasciitilde 10 p.)}
    \begin{itemize}
        \item Reminder of the path we choose with more details: cache / reductions / OP SpMV / ... \textbf{(1 p.)}
        \item Emphasis the comparison between concerned SoA accelerators -> show what is already done and useful, show what is missing
        \item Present the case of banded matrices as an example of structure to leverage \textbf{(1-2 p.)}
        \item SpDCache, architecture presentation with two regions, details on REDs transactions, IBRC vs OBRT \textbf{(7 p.)}
        \item Initial results from Python model (ASAP) and note that a probabilistic model and RTL model will follow
    \end{itemize}
    \item Mathematical analysis of the cache memory accesses for processing sparse matrices \textbf{(\textasciitilde 23 p.)}
    \begin{itemize}
        \item Explain the objectives of this chapter (\textasciitilde formal justification of the concept + cache size/region dimensioning) \textbf{(1 p.)}
        \item Describe the design space of the path we choose (with reminders if necessary) -> introduce base parameters and/or variables, give the exact algorithm... \textbf{(1 p.)}
        \item Show the interest of separate cache regions for A/b and c (formal) -> show that a GC is obviously not efficient when random accesses / confirm why SortCache, PHI, ... exist \textbf{(2 p.)}
        \item Show that A and b are not a limiting factor with a formal reasoning -> give the needed space for GRC when prefetcher (clear statement that gather is not the aim of the thesis) \textbf{(2-3 p.)}
        \item We use a GRC (no modification) for A and b -> get nb loads
        \item Show the interest of a reduction cache region for c -> REDs / reduce the number of mem accesses / link to SoA (see ASAP) \textbf{(2 p.)}
        \item We created a probabilistic model of the reduction cache which is based on the following principles ... (presented in detail in appendix n) / avantage of this model / we implemented this model in C / validation of the model with state of the art article on generic cache formalisation / introduce the graphs for the next sections \textbf{(1-2 p.)}
        \item Show that reduction cache does not handle randomness -> model graph \textbf{(1-2 p.)}
        \item Show RedC 1reg vs RedC 2reg on banded example \textbf{(2-3 p.)}
        \item Show best cache dimensions for the dense part (band) \textbf{(1 p.)}
        \item Show that sparse region (out of band) could use an other region for fine grain density (case of sub bands) -> dimensions \textbf{(1-2 p.)}
        \item Show that sparse region could use adapted HW because of the isolated elements -> dimensions \textbf{(1-2 p.)}
        \item Show how it would work with a uniform matrix -> how does it affects the dimensions \textbf{(1 p.)}
        \item Conclude on our solution and dimensioning \textbf{(1 p.)}
    \end{itemize}
    \item Micro architectiture of SpDCache \textbf{(\textasciitilde 10 p.)}
    \begin{itemize}
        \item Present our hardware solution, based on the concept
        \item Start with basic info on HPDcache (only the one that are useful for SpDCache)
        \item Describe modif in pipeline / hazard management / throughput
        \item Modif in RAM accesses ...
        \item Stats: git / bugs / lines of code / synthesis
    \end{itemize}
    \item RTL simulation results \textbf{(\textasciitilde 10 p.)}
    \begin{itemize}
        \item Experimentations: RTL simulation results -> matrices chosen, generated / metrics / ...
        \item Results / comparison with theory / analysis / conclusion
    \end{itemize}
    \item Future enhancement \textbf{(\textasciitilde 8 p.)}
    \begin{itemize}
        \item Architecture presentation with three regions (add OBRC, CB...) \textbf{(2 p.)}
        \item Architecture of eviction anticipation \textbf{(2 p.)}
        \item Present multicore version -> scale \textbf{(2 p.)}
        \item Present additional optimisation (region config, learn from first iteration) \textbf{(1-2 p.)}
    \end{itemize}
    \item Conclusions \textbf{(\textasciitilde 2 p.)}
    \begin{itemize}
        \item ..
    \end{itemize}
    \item Appendices
    \begin{itemize}
        \item All sparse formats \textbf{(\textasciitilde 5 p.)}
        \item Other accelerators of SoA \textbf{(\textasciitilde 10 p.)}
        \item Probabilistic model of the reduction cache \textbf{(\textasciitilde 10 p.)}
        \item Python model \textbf{(\textasciitilde 1 p.)}
    \end{itemize}
\end{itemize}